% !TEX root =./ECE444_Practice_main.tex
%
% L22 practice set (Antenna Pattern Theory) -- pattern multiplication with a
% patch element factor, scan loss bookkeeping, and reading a measured trace.

\fullwidth{\textbf{Learning Objective 3.6: } Distinguish the array factor from
the true antenna pattern, apply pattern multiplication with a realistic element
pattern, and account for scan loss, coupling, and measurement effects.
\\[4pt]\normalfont\small Show your work. A \textbf{Documentation} line at the end
is required (write \textbf{None} if you did not collaborate).}

\begin{questions}

\question \textbf{The gain budget of the course array.} The ADALM-PHASER
carries $N = 8$ microstrip patch elements spaced $d = 14$ \text{mm} apart. At the
workshop frequency of $10.3\ \ghz$ the wavelength is $29.1$ \text{mm}. Model each
element as an ideal element that captures its share of the aperture, so its
\emph{power} pattern follows the projected-aperture rule
$G_e(\theta) = G_e(0)\cosp{\theta}$ with a broadside element gain of $6.0$ dBi.
\begin{parts}
	\part In one or two sentences, state the difference between the array factor
		and the antenna pattern, and say what the element factor contributes to
		the product.
		\begin{solution}[1.2in]
		The array factor is the contribution of the element positions and their
		relative phases alone; it is what a set of isotropic radiators at those
		positions would produce. The antenna pattern is the array factor
		multiplied by the element factor, which is the pattern of one element in
		the array. The element factor supplies the shape of a single radiator --
		here a broad hemispherical lobe peaked at broadside -- and it sets an
		envelope that the steered array factor cannot exceed.
		\end{solution}

	\begin{minipage}[t]{0.58\linewidth}
		\part Compute the broadside peak gain of the array from the element gain
			and the array gain.
		\begin{solution}[1.1in]
		The eight elements add coherently, so the array gain is
		\eq{ G_{\text{array}} = 10\mylog{8} = 9.0\db }
		and the total broadside gain is
		\eq{ G(0\mydeg) = 6.0\dbi + 9.0\db = 15.0\dbi }
		\end{solution}
	\end{minipage}\hfill%
	\begin{minipage}[t]{0.37\linewidth}
		\raggedleft \vspace{12pt}
		$G(0\mydeg)$ = \ansbox[1.3in]{$15.0\dbi$}
	\end{minipage}

	\begin{minipage}[t]{0.58\linewidth}
		\part The beam is steered to $\theta_0 = 45\mydeg$. Compute the element
			factor penalty in decibels at that angle.
		\begin{solution}[1.1in]
		The projected-aperture rule is stated in power, so the penalty uses
		ten times the logarithm:
		\eq{ 10\mylog{\cosp{45\mydeg}} = 10\mylog{0.707} = -1.5\db }
		\end{solution}
	\end{minipage}\hfill%
	\begin{minipage}[t]{0.37\linewidth}
		\raggedleft \vspace{12pt}
		$\Delta G$ = \ansbox[1.3in]{$-1.5\db$}
	\end{minipage}

	\begin{minipage}[t]{0.58\linewidth}
		\part Compute the peak gain of the array steered to $45\mydeg$.
		\begin{solution}[1.1in]
		Only the element term depends on the steer angle; the array gain of
		$9.0\db$ is unchanged because the elements still add in phase.
		\eq{ G(45\mydeg) &= 6.0\dbi - 1.5\db + 9.0\db \\ &= 13.5\dbi }
		\end{solution}
	\end{minipage}\hfill%
	\begin{minipage}[t]{0.37\linewidth}
		\raggedleft \vspace{12pt}
		$G(45\mydeg)$ = \ansbox[1.3in]{$13.5\dbi$}
	\end{minipage}
\end{parts}

\newpage

\question \textbf{Scan loss and the practical scan limit.} Use the same array
and the same ideal element. Recall from L20 that the array
factor's half-power beamwidth off broadside is
$\theta_{\text{HP}} \approx 0.886\lambda/(Nd\cosp{\theta_0})$, which is
$13.2\mydeg$ at broadside for this array.
\begin{parts}
	\begin{minipage}[t]{0.58\linewidth}
		\part Compute the element factor penalty at $\theta_0 = 30\mydeg$.
		\begin{solution}[1.0in]
		\eq{ 10\mylog{\cosp{30\mydeg}} = 10\mylog{0.866} = -0.6\db }
		\end{solution}
	\end{minipage}\hfill%
	\begin{minipage}[t]{0.37\linewidth}
		\raggedleft \vspace{12pt}
		$\Delta G$ = \ansbox[1.3in]{$-0.6\db$}
	\end{minipage}

	\begin{minipage}[t]{0.58\linewidth}
		\part Compute the element factor penalty at $\theta_0 = 60\mydeg$.
		\begin{solution}[1.0in]
		\eq{ 10\mylog{\cosp{60\mydeg}} = 10\mylog{0.500} = -3.0\db }
		This is the ideal-element bound. A real patch rolls off faster, roughly
		as $\cosp{\theta}^{1.4}$ in power, and measures about $-4.2$ dB at the
		same angle.
		\end{solution}
	\end{minipage}\hfill%
	\begin{minipage}[t]{0.37\linewidth}
		\raggedleft \vspace{12pt}
		$\Delta G$ = \ansbox[1.3in]{$-3.0\db$}
	\end{minipage}

	\begin{minipage}[t]{0.58\linewidth}
		\part Compute the array factor half-power beamwidth at
			$\theta_0 = 60\mydeg$.
		\begin{solution}[1.1in]
		\eq{ \theta_{\text{HP}} = \frac{13.2\mydeg}{\cosp{60\mydeg}}
			= \frac{13.2\mydeg}{0.500} = 26.4\mydeg }
		\end{solution}
	\end{minipage}\hfill%
	\begin{minipage}[t]{0.37\linewidth}
		\raggedleft \vspace{12pt}
		$\theta_{\text{HP}}$ = \ansbox[1.3in]{$26.4\mydeg$}
	\end{minipage}

	\part A radar engineer specifies this array for coverage over
		$\pm 75\mydeg$. Using your numbers from parts (a) through (c), write
		two sentences explaining the trade the engineer is accepting, and state
		why $\pm 60\mydeg$ is the usual limit quoted for patch arrays.
		\begin{solution}[1.4in]
		At $60\mydeg$ the beam has already lost $3\db$ of gain for an ideal element,
		closer to $4\db$ for a real patch, and the array factor's beam has doubled
		in width from $13.2\mydeg$ to $26.4\mydeg$, so the same target returns
		less power and is located less precisely. Pushing to $75\mydeg$ costs
		about $5.9\db$ and widens the beam to nearly $51\mydeg$, so the coverage
		at the edge of the scan is worth much less than the coverage at broadside.
		The $\pm 60\mydeg$ limit is where both penalties are still tolerable;
		beyond it the gain falls faster than the added coverage is worth.
		\end{solution}
\end{parts}

\newpage

\question \textbf{Reading a steered sweep.} A student steers the array to
$\theta_0 = 30\mydeg$ and records the rectangular sweep. The peak reads
$-0.6\db$ relative to the broadside peak and sits at $29.6\mydeg$. The first
sidelobe on the broadside side of the beam reads $-12.2$ dBc and the first
sidelobe on the outboard side reads $-15.2$ dBc. Uniform array factor theory
predicts $-12.8$ dBc for both.
\begin{parts}
	\part Explain in one or two sentences why the measured peak sits at
		$29.6\mydeg$ rather than at the commanded $30\mydeg$.
		\begin{solution}[1.2in]
		The element factor falls monotonically away from broadside, so the
		product of the array factor and the element factor is largest slightly
		inboard of the array factor's own peak. The beam is pulled toward
		broadside. The effect grows with steer angle: a $60\mydeg$ command on
		this array peaks near $56.8\mydeg$.
		\end{solution}

	\part Explain why the two first sidelobes differ by $3$ dB when the
		array factor predicts that they are equal.
		\begin{solution}[1.2in]
		Sidelobe levels are quoted relative to the peak of the same trace, and
		the element factor multiplies each lobe by a different amount. The
		inboard lobe near $7\mydeg$ sits where the element pattern is
		essentially flat, so it is unchanged, while the peak it is referenced to
		fell by $0.6\text{ dB}$; that lobe therefore reads about half a decibel
		\emph{higher} than theory. The outboard lobe near $60\mydeg$ is
		attenuated $3.0\db$ by the element pattern, so it reads well below
		theory.
		\end{solution}

	\begin{minipage}[t]{0.58\linewidth}
		\part Express the inboard sidelobe level relative to the array's
			\emph{broadside} peak rather than relative to the steered peak.
		\begin{solution}[1.1in]
		The steered peak is $0.6\db$ below the broadside peak, so
		\eq{ -12.2\db - 0.6\db = -12.8\db }
		relative to broadside. Referenced this way the lobe matches the array
		factor prediction exactly, which shows that nothing about the sidelobe
		changed -- only the reference did.
		\end{solution}
	\end{minipage}\hfill%
	\begin{minipage}[t]{0.37\linewidth}
		\raggedleft \vspace{12pt}
		Level = \ansbox[1.3in]{$-12.8\db$}
	\end{minipage}

	\part The same student rotates the array by hand and repeats the
		measurement. The new trace has the same shape but every feature is
		shifted $1.8\mydeg$ in the same direction. State whether this is an
		antenna effect or a measurement effect, and give the reason.
		\begin{solution}[1.2in]
		It is a measurement effect. A uniform shift of every feature by the same
		amount, with no change in shape, is a fixture or protractor offset in the
		angle axis. An antenna effect such as beam pulling changes the peak and
		the sidelobes by different amounts and therefore changes the shape of the
		trace, not just its position.
		\end{solution}
\end{parts}

\newpage

\question \textbf{The element and its ground plane.} Take the ideal element
power pattern to be $\cosp{\theta}$ over the front hemisphere and zero behind the
ground plane.
\begin{parts}
	\begin{minipage}[t]{0.58\linewidth}
		\part Find the half-power beamwidth of this element.
		\begin{solution}[1.1in]
		The pattern is stated in power, so half power is where
		\eq{ \cosp{\theta} = 0.5 \quad \rightarrow \quad \theta = 60\mydeg }
		and the full beamwidth is twice that, $120\mydeg$. A real patch is
		steeper and measures nearer $105\mydeg$.
		\end{solution}
	\end{minipage}\hfill%
	\begin{minipage}[t]{0.37\linewidth}
		\raggedleft \vspace{12pt}
		$\theta_{\text{HP}}$ = \ansbox[1.3in]{$120\mydeg$}
	\end{minipage}

	\begin{minipage}[t]{0.58\linewidth}
		\part Compute the directivity of this element, in dBi, integrating over
			the upper hemisphere only.
		\begin{solution}[1.4in]
		The beam solid angle is
		\eq{ \Omega_A &= \int_0^{2\pi}\!\!\int_0^{\pi/2}
			\cosp{\theta} \sin\theta \ d\theta \ d\phi \\
			&= 2\pi \lb \frac{\sin^2\theta}{2} \rb_0^{\pi/2} = \pi }
		so
		\eq{ D_e = \frac{4\pi}{\Omega_A} = \frac{4\pi}{\pi} = 4
			= 6.0\dbi }
		A real patch measures 6.0 to 6.8 dBi, since it is steeper than the ideal
		element, the ground plane is finite, and the element is loaded by its
		neighbors.
		\end{solution}
	\end{minipage}\hfill%
	\begin{minipage}[t]{0.37\linewidth}
		\raggedleft \vspace{12pt}
		$D_e$ = \ansbox[1.3in]{$4$ ($6.0$ dBi)}
	\end{minipage}

	\part A sweep shows a lobe at $-20$ dBc near $\theta = 82\mydeg$. Explain
		what this feature is, and why it cannot be radiation out the back of the
		array.
		\begin{solution}[1.2in]
		The ground plane removes the back hemisphere, and the front-hemisphere
		element pattern falls monotonically toward the horizon, so the antenna
		itself cannot produce a lobe there. The feature is energy arriving from
		a reflection -- a wall, a bench, or a person -- reaching the array from a
		direction the source is not in. Moving the fixture or the reflector
		changes the feature, which is the test that identifies it.
		\end{solution}

	\part The projected-aperture rule is a model. State one angular range where
		you would trust a prediction made with it to about a decibel, and one
		where you would not, and give the reason.
		\begin{solution}[1.2in]
		Inside about $\pm 30\mydeg$ the prediction is good to a few tenths of a
		decibel, because real embedded elements roll off between
		$\cosp{\theta}$ and $\cosp{\theta}^{1.5}$ in power and the choices
		barely differ there. Past about $\pm 60\mydeg$ that same spread is worth
		more than a decibel -- the ideal rule gives $-3.0\db$ where a real patch
		gives about $-4.2\db$ -- and mutual coupling makes each element's
		embedded pattern differ from the average, so a single-element model
		should not be trusted to a decibel at wide scan.
		\end{solution}
\end{parts}

\newpage

\question \textbf{Range, environment, and what a measurement can resolve.} The
course array measures $D = 98$ \text{mm} between the centers of its outer
elements. The lab source is placed $1.0\ \m$ from the array face.
\begin{parts}
	\begin{minipage}[t]{0.58\linewidth}
		\part Compute the far-field distance for this array at $10.3\ \ghz$
			($\lambda = 29.1$ \text{mm}), and state whether the lab range
			qualifies.
		\begin{solution}[1.2in]
		\eq{ r \ge \frac{2D^2}{\lambda}
			= \frac{2\lp0.098\ \m\rp^2}{0.0291\ \m} = 0.66\ \m }
		The $1.0\ \m$ lab range is comfortably beyond this, so the incoming
		wavefront is flat enough that the edge-to-center phase error stays under
		the $22.5\mydeg$ the criterion allows.
		\end{solution}
	\end{minipage}\hfill%
	\begin{minipage}[t]{0.37\linewidth}
		\raggedleft \vspace{12pt}
		$r$ = \ansbox[1.3in]{$0.66\ \m$}
	\end{minipage}

	\begin{minipage}[t]{0.58\linewidth}
		\part A larger array is to be measured on the same $1.0\ \m$ range at
			$10.3\ \ghz$. Find the largest aperture dimension that still satisfies
			the far-field criterion.
		\begin{solution}[1.2in]
		Solve the criterion for $D$:
		\eq{ D = \sqrt{\frac{r\lambda}{2}}
			= \sqrt{\frac{\lp1.0\ \m\rp\lp0.0291\ \m\rp}{2}} = 0.121\ \m }
		so an aperture larger than about $121$ \text{mm} needs a longer range or
		a near-field correction.
		\end{solution}
	\end{minipage}\hfill%
	\begin{minipage}[t]{0.37\linewidth}
		\raggedleft \vspace{12pt}
		$D$ = \ansbox[1.3in]{$121$ \text{mm}}
	\end{minipage}

	\part The GUI sweeps the commanded steer angle in $2.8125\mydeg$ steps.
		Explain how that step size alone can make a measured half-power beamwidth
		differ from theory by one to three degrees.
		\begin{solution}[1.2in]
		The trace exists only at grid points, so each half-power crossing is
		reported at the nearest step rather than at its true angle. Each of the
		two crossings can be off by up to half a step, and the beamwidth is the
		difference between them, so the reported width can be off by roughly a
		full step in either direction before any antenna effect is involved.
		\end{solution}

	\part Name two things you can change about the measurement setup, not the
		antenna, to make the far-out sidelobes of a trace more trustworthy.
		\begin{solution}[1.2in]
		Reduce what the environment contributes: place absorber on the bench and
		on nearby reflecting surfaces, keep people and equipment out of the path,
		and raise the array and source to the same height away from large flat
		surfaces. Repeat the measurement with the whole fixture moved or rotated;
		features that move with the room are environmental, and features that
		stay with the array are the antenna. Increasing the source power raises
		the trace above the sweep noise floor, which is what limits any reading
		more than about 23 dB below the peak.
		\end{solution}
\end{parts}

\vspace{1.5em}
\noindent\textbf{Documentation:} \rule[-2pt]{0.6\linewidth}{0.4pt}

\end{questions}
